% ============================================================================
% OLMo in-family self-scoring matrix (grader-robustness appendix).
% Every number in this section resolves through generated artifacts:
%   results/rlhf_experiment/paper_macros_matrix.tex      (inline macros)
%   results/rlhf_experiment/tables_matrix/*.tex          (table bodies)
% both emitted by scripts/rlhf_experiment/7_matrix_paper_outputs.py, which
% consumes matrix_summary.json from 6_matrix_analysis.py.
% Raw findings write-up: results/rlhf_experiment/matrix/FINDINGS.md
% ============================================================================

% Load generated macros BEFORE the section body so prose references resolve.
\IfFileExists{../results/rlhf_experiment/paper_macros_matrix.tex}{%
  % Auto-generated by scripts/rlhf_experiment/7_matrix_paper_outputs.py -- do not hand-edit.
% Inline-number macros for the OLMo self-scoring matrix experiment
% (paper/sections/appG_olmo_self_scoring.tex and FINDINGS.md prose).
\newcommand{\olmoMxAlpacaN}{150}
\newcommand{\olmoMxAlpacaDThetaQwenGenBase}{0.481}
\newcommand{\olmoMxAlpacaCThetaQwenGenBase}{0.436}
\newcommand{\olmoMxAlpacaAnThetaQwenGenBase}{1.142}
\newcommand{\olmoMxAlpacaDThetaQwenGenSft}{0.329}
\newcommand{\olmoMxAlpacaCThetaQwenGenSft}{0.674}
\newcommand{\olmoMxAlpacaAnThetaQwenGenSft}{0.512}
\newcommand{\olmoMxAlpacaDThetaQwenGenDpo}{0.286}
\newcommand{\olmoMxAlpacaCThetaQwenGenDpo}{0.694}
\newcommand{\olmoMxAlpacaAnThetaQwenGenDpo}{0.427}
\newcommand{\olmoMxAlpacaDThetaQwenGenInstruct}{0.281}
\newcommand{\olmoMxAlpacaCThetaQwenGenInstruct}{0.698}
\newcommand{\olmoMxAlpacaAnThetaQwenGenInstruct}{0.419}
\newcommand{\olmoMxAlpacaAnFoldThetaQwen}{2.7}
\newcommand{\olmoMxAlpacaDThetaBaseGenBase}{0.494}
\newcommand{\olmoMxAlpacaCThetaBaseGenBase}{0.420}
\newcommand{\olmoMxAlpacaAnThetaBaseGenBase}{1.224}
\newcommand{\olmoMxAlpacaDThetaBaseGenSft}{0.345}
\newcommand{\olmoMxAlpacaCThetaBaseGenSft}{0.607}
\newcommand{\olmoMxAlpacaAnThetaBaseGenSft}{0.603}
\newcommand{\olmoMxAlpacaDThetaBaseGenDpo}{0.298}
\newcommand{\olmoMxAlpacaCThetaBaseGenDpo}{0.603}
\newcommand{\olmoMxAlpacaAnThetaBaseGenDpo}{0.518}
\newcommand{\olmoMxAlpacaDThetaBaseGenInstruct}{0.292}
\newcommand{\olmoMxAlpacaCThetaBaseGenInstruct}{0.605}
\newcommand{\olmoMxAlpacaAnThetaBaseGenInstruct}{0.511}
\newcommand{\olmoMxAlpacaAnFoldThetaBase}{2.4}
\newcommand{\olmoMxAlpacaDThetaSftGenBase}{0.487}
\newcommand{\olmoMxAlpacaCThetaSftGenBase}{0.395}
\newcommand{\olmoMxAlpacaAnThetaSftGenBase}{1.288}
\newcommand{\olmoMxAlpacaDThetaSftGenSft}{0.356}
\newcommand{\olmoMxAlpacaCThetaSftGenSft}{0.642}
\newcommand{\olmoMxAlpacaAnThetaSftGenSft}{0.591}
\newcommand{\olmoMxAlpacaDThetaSftGenDpo}{0.312}
\newcommand{\olmoMxAlpacaCThetaSftGenDpo}{0.648}
\newcommand{\olmoMxAlpacaAnThetaSftGenDpo}{0.503}
\newcommand{\olmoMxAlpacaDThetaSftGenInstruct}{0.307}
\newcommand{\olmoMxAlpacaCThetaSftGenInstruct}{0.650}
\newcommand{\olmoMxAlpacaAnThetaSftGenInstruct}{0.497}
\newcommand{\olmoMxAlpacaAnFoldThetaSft}{2.6}
\newcommand{\olmoMxAlpacaDThetaDpoGenBase}{0.494}
\newcommand{\olmoMxAlpacaCThetaDpoGenBase}{0.341}
\newcommand{\olmoMxAlpacaAnThetaDpoGenBase}{1.530}
\newcommand{\olmoMxAlpacaDThetaDpoGenSft}{0.398}
\newcommand{\olmoMxAlpacaCThetaDpoGenSft}{0.615}
\newcommand{\olmoMxAlpacaAnThetaDpoGenSft}{0.696}
\newcommand{\olmoMxAlpacaDThetaDpoGenDpo}{0.354}
\newcommand{\olmoMxAlpacaCThetaDpoGenDpo}{0.646}
\newcommand{\olmoMxAlpacaAnThetaDpoGenDpo}{0.574}
\newcommand{\olmoMxAlpacaDThetaDpoGenInstruct}{0.349}
\newcommand{\olmoMxAlpacaCThetaDpoGenInstruct}{0.649}
\newcommand{\olmoMxAlpacaAnThetaDpoGenInstruct}{0.566}
\newcommand{\olmoMxAlpacaAnFoldThetaDpo}{2.7}
\newcommand{\olmoMxAlpacaDThetaInstructGenBase}{0.494}
\newcommand{\olmoMxAlpacaCThetaInstructGenBase}{0.331}
\newcommand{\olmoMxAlpacaAnThetaInstructGenBase}{1.579}
\newcommand{\olmoMxAlpacaDThetaInstructGenSft}{0.404}
\newcommand{\olmoMxAlpacaCThetaInstructGenSft}{0.608}
\newcommand{\olmoMxAlpacaAnThetaInstructGenSft}{0.715}
\newcommand{\olmoMxAlpacaDThetaInstructGenDpo}{0.359}
\newcommand{\olmoMxAlpacaCThetaInstructGenDpo}{0.641}
\newcommand{\olmoMxAlpacaAnThetaInstructGenDpo}{0.586}
\newcommand{\olmoMxAlpacaDThetaInstructGenInstruct}{0.354}
\newcommand{\olmoMxAlpacaCThetaInstructGenInstruct}{0.645}
\newcommand{\olmoMxAlpacaAnThetaInstructGenInstruct}{0.578}
\newcommand{\olmoMxAlpacaAnFoldThetaInstruct}{2.7}
\newcommand{\olmoMxAlpacaPholmThetaQwenBaseSft}{1.0 \times 10^{-23}}
\newcommand{\olmoMxAlpacaPholmThetaQwenSftDpo}{2.8 \times 10^{-13}}
\newcommand{\olmoMxAlpacaPholmThetaBaseBaseSft}{2.5 \times 10^{-24}}
\newcommand{\olmoMxAlpacaPholmThetaBaseSftDpo}{3.8 \times 10^{-12}}
\newcommand{\olmoMxAlpacaPholmThetaSftBaseSft}{6.4 \times 10^{-22}}
\newcommand{\olmoMxAlpacaPholmThetaSftSftDpo}{6.3 \times 10^{-10}}
\newcommand{\olmoMxAlpacaPholmThetaDpoBaseSft}{1.2 \times 10^{-16}}
\newcommand{\olmoMxAlpacaPholmThetaDpoSftDpo}{1.3 \times 10^{-8}}
\newcommand{\olmoMxAlpacaPholmThetaInstructBaseSft}{8.0 \times 10^{-16}}
\newcommand{\olmoMxAlpacaPholmThetaInstructSftDpo}{1.3 \times 10^{-8}}
\newcommand{\olmoMxAlpacaRThetaQwen}{1.68}
\newcommand{\olmoMxAlpacaRLoThetaQwen}{1.61}
\newcommand{\olmoMxAlpacaRHiThetaQwen}{1.76}
\newcommand{\olmoMxAlpacaRThetaBase}{1.66}
\newcommand{\olmoMxAlpacaRLoThetaBase}{1.60}
\newcommand{\olmoMxAlpacaRHiThetaBase}{1.72}
\newcommand{\olmoMxAlpacaRThetaSft}{1.56}
\newcommand{\olmoMxAlpacaRLoThetaSft}{1.50}
\newcommand{\olmoMxAlpacaRHiThetaSft}{1.62}
\newcommand{\olmoMxAlpacaRThetaDpo}{1.40}
\newcommand{\olmoMxAlpacaRLoThetaDpo}{1.34}
\newcommand{\olmoMxAlpacaRHiThetaDpo}{1.45}
\newcommand{\olmoMxAlpacaRThetaInstruct}{1.38}
\newcommand{\olmoMxAlpacaRLoThetaInstruct}{1.33}
\newcommand{\olmoMxAlpacaRHiThetaInstruct}{1.43}
\newcommand{\olmoMxAlpacaVarStagePct}{88.0}
\newcommand{\olmoMxAlpacaVarStageInFamPct}{89.8}
\newcommand{\olmoMxAlpacaVarWithinStagePct}{86.7}
\newcommand{\olmoMxAlpacaVarGraderPct}{9.6}
\newcommand{\olmoMxAlpacaVarGraderInFamPct}{7.9}
\newcommand{\olmoMxAlpacaVarWithinGraderPct}{9.3}
\newcommand{\olmoMxAlpacaVarInteractionPct}{2.3}
\newcommand{\olmoMxAlpacaVarInteractionInFamPct}{2.3}
\newcommand{\olmoMxAlpacaVarWithinInteractionPct}{4.0}
\newcommand{\olmoMxAlpacaVarPromptPct}{23.8}
\newcommand{\olmoMxNbN}{39}
\newcommand{\olmoMxNbDThetaQwenGenBase}{0.481}
\newcommand{\olmoMxNbCThetaQwenGenBase}{0.391}
\newcommand{\olmoMxNbAnThetaQwenGenBase}{1.257}
\newcommand{\olmoMxNbDThetaQwenGenSft}{0.369}
\newcommand{\olmoMxNbCThetaQwenGenSft}{0.605}
\newcommand{\olmoMxNbAnThetaQwenGenSft}{0.642}
\newcommand{\olmoMxNbDThetaQwenGenDpo}{0.312}
\newcommand{\olmoMxNbCThetaQwenGenDpo}{0.637}
\newcommand{\olmoMxNbAnThetaQwenGenDpo}{0.518}
\newcommand{\olmoMxNbDThetaQwenGenInstruct}{0.303}
\newcommand{\olmoMxNbCThetaQwenGenInstruct}{0.639}
\newcommand{\olmoMxNbAnThetaQwenGenInstruct}{0.504}
\newcommand{\olmoMxNbAnFoldThetaQwen}{2.5}
\newcommand{\olmoMxNbDThetaBaseGenBase}{0.492}
\newcommand{\olmoMxNbCThetaBaseGenBase}{0.374}
\newcommand{\olmoMxNbAnThetaBaseGenBase}{1.353}
\newcommand{\olmoMxNbDThetaBaseGenSft}{0.358}
\newcommand{\olmoMxNbCThetaBaseGenSft}{0.554}
\newcommand{\olmoMxNbAnThetaBaseGenSft}{0.677}
\newcommand{\olmoMxNbDThetaBaseGenDpo}{0.302}
\newcommand{\olmoMxNbCThetaBaseGenDpo}{0.553}
\newcommand{\olmoMxNbAnThetaBaseGenDpo}{0.584}
\newcommand{\olmoMxNbDThetaBaseGenInstruct}{0.290}
\newcommand{\olmoMxNbCThetaBaseGenInstruct}{0.551}
\newcommand{\olmoMxNbAnThetaBaseGenInstruct}{0.567}
\newcommand{\olmoMxNbAnFoldThetaBase}{2.4}
\newcommand{\olmoMxNbDThetaSftGenBase}{0.483}
\newcommand{\olmoMxNbCThetaSftGenBase}{0.352}
\newcommand{\olmoMxNbAnThetaSftGenBase}{1.417}
\newcommand{\olmoMxNbDThetaSftGenSft}{0.374}
\newcommand{\olmoMxNbCThetaSftGenSft}{0.599}
\newcommand{\olmoMxNbAnThetaSftGenSft}{0.662}
\newcommand{\olmoMxNbDThetaSftGenDpo}{0.316}
\newcommand{\olmoMxNbCThetaSftGenDpo}{0.609}
\newcommand{\olmoMxNbAnThetaSftGenDpo}{0.557}
\newcommand{\olmoMxNbDThetaSftGenInstruct}{0.306}
\newcommand{\olmoMxNbCThetaSftGenInstruct}{0.608}
\newcommand{\olmoMxNbAnThetaSftGenInstruct}{0.541}
\newcommand{\olmoMxNbAnFoldThetaSft}{2.6}
\newcommand{\olmoMxNbDThetaDpoGenBase}{0.494}
\newcommand{\olmoMxNbCThetaDpoGenBase}{0.305}
\newcommand{\olmoMxNbAnThetaDpoGenBase}{1.679}
\newcommand{\olmoMxNbDThetaDpoGenSft}{0.413}
\newcommand{\olmoMxNbCThetaDpoGenSft}{0.570}
\newcommand{\olmoMxNbAnThetaDpoGenSft}{0.779}
\newcommand{\olmoMxNbDThetaDpoGenDpo}{0.357}
\newcommand{\olmoMxNbCThetaDpoGenDpo}{0.606}
\newcommand{\olmoMxNbAnThetaDpoGenDpo}{0.636}
\newcommand{\olmoMxNbDThetaDpoGenInstruct}{0.351}
\newcommand{\olmoMxNbCThetaDpoGenInstruct}{0.608}
\newcommand{\olmoMxNbAnThetaDpoGenInstruct}{0.623}
\newcommand{\olmoMxNbAnFoldThetaDpo}{2.7}
\newcommand{\olmoMxNbDThetaInstructGenBase}{0.495}
\newcommand{\olmoMxNbCThetaInstructGenBase}{0.296}
\newcommand{\olmoMxNbAnThetaInstructGenBase}{1.735}
\newcommand{\olmoMxNbDThetaInstructGenSft}{0.419}
\newcommand{\olmoMxNbCThetaInstructGenSft}{0.563}
\newcommand{\olmoMxNbAnThetaInstructGenSft}{0.801}
\newcommand{\olmoMxNbDThetaInstructGenDpo}{0.362}
\newcommand{\olmoMxNbCThetaInstructGenDpo}{0.600}
\newcommand{\olmoMxNbAnThetaInstructGenDpo}{0.651}
\newcommand{\olmoMxNbDThetaInstructGenInstruct}{0.356}
\newcommand{\olmoMxNbCThetaInstructGenInstruct}{0.603}
\newcommand{\olmoMxNbAnThetaInstructGenInstruct}{0.637}
\newcommand{\olmoMxNbAnFoldThetaInstruct}{2.7}
\newcommand{\olmoMxNbPholmThetaQwenBaseSft}{1.2 \times 10^{-7}}
\newcommand{\olmoMxNbPholmThetaQwenSftDpo}{6.6 \times 10^{-6}}
\newcommand{\olmoMxNbPholmThetaBaseBaseSft}{3.1 \times 10^{-9}}
\newcommand{\olmoMxNbPholmThetaBaseSftDpo}{1.1 \times 10^{-5}}
\newcommand{\olmoMxNbPholmThetaSftBaseSft}{2.2 \times 10^{-7}}
\newcommand{\olmoMxNbPholmThetaSftSftDpo}{1.1 \times 10^{-5}}
\newcommand{\olmoMxNbPholmThetaDpoBaseSft}{4.4 \times 10^{-5}}
\newcommand{\olmoMxNbPholmThetaDpoSftDpo}{4.4 \times 10^{-5}}
\newcommand{\olmoMxNbPholmThetaInstructBaseSft}{4.4 \times 10^{-5}}
\newcommand{\olmoMxNbPholmThetaInstructSftDpo}{4.4 \times 10^{-5}}
\newcommand{\olmoMxNbRThetaQwen}{1.54}
\newcommand{\olmoMxNbRLoThetaQwen}{1.41}
\newcommand{\olmoMxNbRHiThetaQwen}{1.68}
\newcommand{\olmoMxNbRThetaBase}{1.63}
\newcommand{\olmoMxNbRLoThetaBase}{1.51}
\newcommand{\olmoMxNbRHiThetaBase}{1.76}
\newcommand{\olmoMxNbRThetaSft}{1.53}
\newcommand{\olmoMxNbRLoThetaSft}{1.41}
\newcommand{\olmoMxNbRHiThetaSft}{1.66}
\newcommand{\olmoMxNbRThetaDpo}{1.38}
\newcommand{\olmoMxNbRLoThetaDpo}{1.28}
\newcommand{\olmoMxNbRHiThetaDpo}{1.49}
\newcommand{\olmoMxNbRThetaInstruct}{1.37}
\newcommand{\olmoMxNbRLoThetaInstruct}{1.27}
\newcommand{\olmoMxNbRHiThetaInstruct}{1.47}
\newcommand{\olmoMxNbVarStagePct}{90.1}
\newcommand{\olmoMxNbVarStageInFamPct}{89.1}
\newcommand{\olmoMxNbVarWithinStagePct}{88.4}
\newcommand{\olmoMxNbVarGraderPct}{8.0}
\newcommand{\olmoMxNbVarGraderInFamPct}{8.7}
\newcommand{\olmoMxNbVarWithinGraderPct}{7.5}
\newcommand{\olmoMxNbVarInteractionPct}{1.9}
\newcommand{\olmoMxNbVarInteractionInFamPct}{2.3}
\newcommand{\olmoMxNbVarWithinInteractionPct}{4.1}
\newcommand{\olmoMxNbVarPromptPct}{32.8}
\newcommand{\olmoMxNContrasts}{20}
\newcommand{\olmoMxNContrastsSig}{20}
\newcommand{\olmoMxMaxPholm}{4.4 \times 10^{-5}}
\newcommand{\olmoMxMinR}{1.37}
\newcommand{\olmoMxMinRLo}{1.27}
%
}{}

\section{Grader Robustness: The OLMo Self-Scoring Matrix}\label{app:olmo-self-scoring}

The RLHF case study (Section~\ref{sec:rlhf-experiment}) scores all four OLMo-2-7B pipeline stages with a single cross-family grader, Qwen2.5-3B.
A natural objection is that the measured diversity drop could be a scoring artifact: if $D_{Ca_n}$ partly rewards responses that resemble the grader $\theta$, then a $\theta$ far from the aligned OLMo checkpoints would depress their scores regardless of their actual diversity.
This appendix stress-tests that concern by re-scoring the same fixed generations with every OLMo pipeline stage as the grader, including the configuration that maximally advantages the aligned stages: the RLVR-tuned Instruct model grading its own family's outputs, in deliberate violation of the base-model-$\theta$ requirement of Section~\ref{sec:rlhf-experiment}.

\paragraph{Design.}
We hold the length-matched generations, the ${K{=}10}$ responses per prompt, and the 25 permutations (seed 42) of Section~\ref{sec:rlhf-experiment} fixed, and vary only $\theta$: the four released OLMo-2-1124-7B stages (base, SFT, DPO, Instruct) plus the published cross-family Qwen2.5-3B row, giving a $5 \times 4$ grader-by-generator matrix of mean $D_{Ca_n}$ on each prompt set (\olmoMxAlpacaN{} AlpacaEval and \olmoMxNbN{} NoveltyBench-curated prompts).
The confirmatory prediction in every grader row is the adjacent chain base $>$ SFT $>$ DPO, tested by one-sided paired Wilcoxon signed-rank tests with Holm correction within each prompt set; the DPO-vs-Instruct step carries no directional prediction (Section~\ref{sec:rlhf-experiment}) and is excluded from the confirmatory family.

\begin{table}[h]
\centering
\caption{Mean per-prompt $D_{Ca_n}$ for every grader $\theta$ (rows) scoring every OLMo-2-7B generation stage (columns), on both length-matched prompt sets. The first row is the cross-family configuration reported in Section~\ref{sec:rlhf-experiment}; the remaining rows are in-family graders, including each aligned stage scoring its own outputs on the diagonal.}
\label{tab:olmo-matrix-D}
\small
\IfFileExists{../results/rlhf_experiment/tables_matrix/olmo_matrix_D.tex}{%
  % Auto-generated by scripts/rlhf_experiment/7_matrix_paper_outputs.py -- do not hand-edit.
\begin{tabular}{lrrrr}
\toprule
\multicolumn{5}{l}{\textbf{AlpacaEval} ($n=150$ prompts)} \\
\midrule
$\theta$ (grader) & Base & SFT & DPO & Instruct (RLVR) \\
\midrule
Qwen2.5-3B (cross-family) & 0.481 & 0.329 & 0.286 & 0.281 \\
OLMo base & 0.494 & 0.345 & 0.298 & 0.292 \\
OLMo SFT & 0.487 & 0.356 & 0.312 & 0.307 \\
OLMo DPO & 0.494 & 0.398 & 0.354 & 0.349 \\
OLMo Instruct (RLVR) & 0.494 & 0.404 & 0.359 & 0.354 \\
\addlinespace
Llama-3.1-8B (cross-family) & 0.484 & 0.312 & 0.263 & 0.258 \\
GPT-2 124M (cross-family) ($n{=}137$) & 0.492 & 0.388 & 0.333 & 0.324 \\
\midrule
\multicolumn{5}{l}{\textbf{NoveltyBench curated} ($n=39$ prompts)} \\
\midrule
$\theta$ (grader) & Base & SFT & DPO & Instruct (RLVR) \\
\midrule
Qwen2.5-3B (cross-family) & 0.481 & 0.369 & 0.312 & 0.303 \\
OLMo base & 0.492 & 0.358 & 0.302 & 0.290 \\
OLMo SFT & 0.483 & 0.374 & 0.316 & 0.306 \\
OLMo DPO & 0.494 & 0.413 & 0.357 & 0.351 \\
OLMo Instruct (RLVR) & 0.495 & 0.419 & 0.362 & 0.356 \\
\addlinespace
Llama-3.1-8B (cross-family) & 0.489 & 0.339 & 0.278 & 0.267 \\
GPT-2 124M (cross-family) & 0.496 & 0.401 & 0.336 & 0.322 \\
\bottomrule
\end{tabular}
%
}{%
  \texttt{[table appears after 7\_matrix\_paper\_outputs.py runs]}%
}
\end{table}

\paragraph{The drop is grader-robust.}
Table~\ref{tab:olmo-matrix-D} shows that $D_{Ca_n}$ decreases along base $>$ SFT $>$ DPO in every grader row on both prompt sets, and all \olmoMxNContrastsSig{} of \olmoMxNContrasts{} confirmatory contrasts are significant after Holm correction (largest adjusted $p = \olmoMxMaxPholm$).
In the adversarial row, OLMo-Instruct grading its own family, the ordering still holds: base scores $\olmoMxAlpacaDThetaInstructGenBase$ against $\olmoMxAlpacaDThetaInstructGenInstruct$ for the Instruct stage's own outputs on AlpacaEval (base $>$ SFT at $p_{\mathrm{Holm}} = \olmoMxAlpacaPholmThetaInstructBaseSft$, SFT $>$ DPO at $p_{\mathrm{Holm}} = \olmoMxAlpacaPholmThetaInstructSftDpo$).
The fold-change $R_\theta = \bar{D}_{\mathrm{base}} / \bar{D}_{\mathrm{DPO}}$ (Table~\ref{tab:olmo-matrix-R}) shrinks as the grader itself becomes more aligned, from $\olmoMxAlpacaRThetaQwen$ under Qwen2.5-3B to $\olmoMxAlpacaRThetaInstruct$ under OLMo-Instruct on AlpacaEval, but never approaches parity: no point estimate falls below \olmoMxMinR{} and no bootstrap 95\% CI lower bound falls below \olmoMxMinRLo{}.

\begin{table}[h]
\centering
\caption{Fold-change $R_\theta = \bar{D}_{\mathrm{base}} / \bar{D}_{\mathrm{DPO}}$ per grader row with bootstrap 95\% CIs (10{,}000 resamples over prompts). More-aligned graders see a smaller but still clearly above-parity drop.}
\label{tab:olmo-matrix-R}
\small
\IfFileExists{../results/rlhf_experiment/tables_matrix/olmo_matrix_R.tex}{%
  % Auto-generated by scripts/rlhf_experiment/7_matrix_paper_outputs.py -- do not hand-edit.
\begin{tabular}{lcccc}
\toprule
 & \multicolumn{2}{c}{AlpacaEval} & \multicolumn{2}{c}{NB curated} \\
\cmidrule(lr){2-3}\cmidrule(lr){4-5}
$\theta$ (grader) & $R_\theta$ & 95\% CI & $R_\theta$ & 95\% CI \\
\midrule
Qwen2.5-3B (cross-family) & 1.68 & [1.61, 1.76] & 1.54 & [1.41, 1.68] \\
OLMo base & 1.66 & [1.60, 1.72] & 1.63 & [1.51, 1.76] \\
OLMo SFT & 1.56 & [1.50, 1.62] & 1.53 & [1.41, 1.66] \\
OLMo DPO & 1.40 & [1.34, 1.45] & 1.38 & [1.28, 1.49] \\
OLMo Instruct (RLVR) & 1.38 & [1.33, 1.43] & 1.37 & [1.27, 1.47] \\
\bottomrule
\end{tabular}
%
}{%
  \texttt{[table appears after 7\_matrix\_paper\_outputs.py runs]}%
}
\end{table}

\paragraph{Variance decomposition: who is graded matters far more than who grades.}
To quantify how much of the variation in $D_{Ca_n}$ is attributable to grader identity versus the model being graded, we exploit the fully crossed, balanced design and partition sums of squares in two complementary ways.
At the cell-mean level, we decompose the $5 \times 4$ matrix of prompt-averaged $D_{Ca_n}$ into a grader main effect, a generator-stage main effect, and their interaction (the residual of the additive fit).
At the per-prompt level, we decompose the full grader $\times$ stage $\times$ prompt tensor; prompt identity is a nuisance factor (absorbing \olmoMxAlpacaVarPromptPct\% and \olmoMxNbVarPromptPct\% of total sum of squares on the two prompt sets), so we report shares of the within-prompt sum of squares, crediting each factor with its main effect plus its interaction with prompt.
Table~\ref{tab:olmo-matrix-variance} shows the result: the generator stage accounts for \olmoMxAlpacaVarStagePct\% (AlpacaEval) and \olmoMxNbVarStagePct\% (NB-curated) of cell-mean variance, grader identity for only \olmoMxAlpacaVarGraderPct\% and \olmoMxNbVarGraderPct\%, and the interaction for \olmoMxAlpacaVarInteractionPct\% and \olmoMxNbVarInteractionPct\%; the within-prompt shares are nearly identical.
Restricting to the in-family $4 \times 4$ submatrix (dropping the cross-family Qwen row) barely moves the split (stage \olmoMxAlpacaVarStageInFamPct\% / grader \olmoMxAlpacaVarGraderInFamPct\% on AlpacaEval; stage \olmoMxNbVarStageInFamPct\% / grader \olmoMxNbVarGraderInFamPct\% on NB-curated), so the small grader share is not driven by the family switch.
$D_{Ca_n}$ is therefore overwhelmingly a property of the policy being measured, not of the grader measuring it.

\begin{table}[h]
\centering
\caption{Share of variance in $D_{Ca_n}$ attributable to the generator stage, the grader $\theta$, and their interaction, for all five grader rows. Cell-mean columns: two-way sum-of-squares decomposition of the $5 \times 4$ matrix in Table~\ref{tab:olmo-matrix-D}. Within-prompt columns: three-way decomposition of the grader $\times$ stage $\times$ prompt tensor, reported as shares of the sum of squares remaining after removing the prompt main effect, with each factor credited with its interaction with prompt.}
\label{tab:olmo-matrix-variance}
\small
\IfFileExists{../results/rlhf_experiment/tables_matrix/olmo_matrix_variance.tex}{%
  % Auto-generated by scripts/rlhf_experiment/7_matrix_paper_outputs.py -- do not hand-edit.
\begin{tabular}{lcccc}
\toprule
 & \multicolumn{2}{c}{AlpacaEval} & \multicolumn{2}{c}{NB curated} \\
\cmidrule(lr){2-3}\cmidrule(lr){4-5}
Variance source & cell means & within-prompt & cell means & within-prompt \\
\midrule
\multicolumn{5}{l}{\textbf{Pre-registered graders (5)}} \\
\midrule
Generator stage & 88.0\% & 86.7\% & 90.1\% & 88.4\% \\
Grader $\theta$ & 9.6\% & 9.3\% & 8.0\% & 7.5\% \\
Grader $\times$ stage interaction & 2.3\% & 4.0\% & 1.9\% & 4.1\% \\
\midrule
\multicolumn{5}{l}{\textbf{Extended graders (7, common prompts)}} \\
\midrule
Generator stage & 86.8\% & 82.8\% & 89.0\% & 82.3\% \\
Grader $\theta$ & 10.6\% & 11.7\% & 8.7\% & 12.0\% \\
Grader $\times$ stage interaction & 2.6\% & 5.5\% & 2.3\% & 5.7\% \\
\bottomrule
\end{tabular}
%
}{%
  \texttt{[table appears after 7\_matrix\_paper\_outputs.py runs]}%
}
\end{table}

\paragraph{Mechanism.}
The drop is carried by the mutual-predictability term $a_n$, not by coherence: across base $\to$ Instruct generations under the Qwen grader, mean $a_n$ falls \olmoMxAlpacaAnFoldThetaQwen$\times$ (from \olmoMxAlpacaAnThetaQwenGenBase{} to \olmoMxAlpacaAnThetaQwenGenInstruct{} bits/byte on AlpacaEval) while mean $C$ \emph{rises} (from \olmoMxAlpacaCThetaQwenGenBase{} to \olmoMxAlpacaCThetaQwenGenInstruct), and the same pattern holds in every grader row.
The only self-similarity signal in the matrix points against a spurious drop: on base-stage text, coherence is highest under the base grader and declines as the grader becomes more aligned (from \olmoMxAlpacaCThetaBaseGenBase{} to \olmoMxAlpacaCThetaInstructGenBase{} on AlpacaEval), and aligned graders assign base-stage text a \emph{higher} $a_n$ than the base grader does (\olmoMxAlpacaAnThetaBaseGenBase{} vs.\ \olmoMxAlpacaAnThetaInstructGenBase{} bits/byte), widening rather than narrowing the measured gap.
We retain the base-model-$\theta$ requirement as a design principle, since an aligned $\theta$ conflates coherence-as-fluency with coherence-as-alignment, but the post-training diversity drop of Section~\ref{sec:rlhf-experiment} does not depend on it.

% --- macro fallbacks (only fire if the generated files are absent) ---
\providecommand{\olmoMxAlpacaN}{??}
\providecommand{\olmoMxNbN}{??}
\providecommand{\olmoMxNContrasts}{??}
\providecommand{\olmoMxNContrastsSig}{??}
\providecommand{\olmoMxMaxPholm}{??}
\providecommand{\olmoMxMinR}{??}
\providecommand{\olmoMxMinRLo}{??}
\providecommand{\olmoMxAlpacaDThetaInstructGenBase}{??}
\providecommand{\olmoMxAlpacaDThetaInstructGenInstruct}{??}
\providecommand{\olmoMxAlpacaPholmThetaInstructBaseSft}{??}
\providecommand{\olmoMxAlpacaPholmThetaInstructSftDpo}{??}
\providecommand{\olmoMxAlpacaRThetaQwen}{??}
\providecommand{\olmoMxAlpacaRThetaInstruct}{??}
\providecommand{\olmoMxAlpacaVarPromptPct}{??}
\providecommand{\olmoMxNbVarPromptPct}{??}
\providecommand{\olmoMxAlpacaVarStagePct}{??}
\providecommand{\olmoMxNbVarStagePct}{??}
\providecommand{\olmoMxAlpacaVarGraderPct}{??}
\providecommand{\olmoMxNbVarGraderPct}{??}
\providecommand{\olmoMxAlpacaVarInteractionPct}{??}
\providecommand{\olmoMxNbVarInteractionPct}{??}
\providecommand{\olmoMxAlpacaVarStageInFamPct}{??}
\providecommand{\olmoMxAlpacaVarGraderInFamPct}{??}
\providecommand{\olmoMxNbVarStageInFamPct}{??}
\providecommand{\olmoMxNbVarGraderInFamPct}{??}
\providecommand{\olmoMxAlpacaAnFoldThetaQwen}{??}
\providecommand{\olmoMxAlpacaAnThetaQwenGenBase}{??}
\providecommand{\olmoMxAlpacaAnThetaQwenGenInstruct}{??}
\providecommand{\olmoMxAlpacaCThetaQwenGenBase}{??}
\providecommand{\olmoMxAlpacaCThetaQwenGenInstruct}{??}
\providecommand{\olmoMxAlpacaCThetaBaseGenBase}{??}
\providecommand{\olmoMxAlpacaCThetaInstructGenBase}{??}
\providecommand{\olmoMxAlpacaAnThetaBaseGenBase}{??}
\providecommand{\olmoMxAlpacaAnThetaInstructGenBase}{??}
